\section{RT-2, RT-X, OpenVLA y Action Head}

Ya se explicó que los modelos base para robots necesitan datos de robots reales y políticas que funcionen entre tareas; esta sección avanza hacia el grupo de artículos más cercano a la línea principal actual de VLA. Aquí la pregunta cambia: si un VLM ya aprendió una gran cantidad de conocimiento visual y lingüístico a partir de imágenes y textos de internet, ¿puede el robot transferir ese conocimiento a la salida de acciones, en lugar de volver a aprender el sentido común sobre el mundo a partir de una pequeña cantidad de datos de robots?

Esta sección entra en la línea principal actual de VLA. RT-2 se considera uno de los trabajos pioneros de VLA: transfiere el conocimiento visual y lingüístico a escala web al control de robots. La idea central es esta: un VLM preentrenado ya domina una gran cantidad de conocimiento visual y lingüístico; ¿es posible, mediante datos de acciones de robots, hacer que produzca directamente acciones (action)? Esto conecta el conocimiento de internet con el control de robots.

\begin{figure}[H]
\centering
\includegraphics[width=0.96\textwidth]{ch23-05001.jpg}
\caption{RT-2: Vision-Language-Action Models transfer web knowledge to robotic control. Fotograma de la pantalla proyectada, aprox. 01:23:21.}
\end{figure}

\begin{knowledgebox}{Lectura de la figura: RT-2 usa un VLM como backbone para producir acciones del robot}
En la figura se observa que RT-2 conecta el VLM con las acciones (action) del robot. El núcleo es la transferencia: el modelo no aprende desde cero solo con datos de robots, sino que transfiere el conocimiento visual y lingüístico de internet al control de robots.
\end{knowledgebox}

\subsection{RT-X y Open X-Embodiment}

Esta subsección desplaza el foco de la estructura de un modelo individual a la infraestructura de datos. RT-2 demostró que transferir un VLM al control de robots tiene valor, pero si los datos de entrenamiento provienen solo de unos pocos robots y unas pocas tareas, el modelo sigue quedando fácilmente atado a la plataforma robótica, al escenario y al espacio de acciones. Por eso, la pregunta central de RT-X/Open X-Embodiment es cómo reunir la experiencia dispersa en distintos laboratorios y distintos robots en datos generales con los que se pueda entrenar.

RT-X y Open X-Embodiment ponen el énfasis en el conjunto de datos y en los modelos que funcionan entre plataformas robóticas. Los datos de robots son escasos y fragmentados, y a un solo laboratorio le resulta difícil cubrir suficientes tareas y plataformas robóticas. Open X-Embodiment reúne datos de múltiples robots y múltiples tareas para intentar entrenar modelos que funcionen entre plataformas robóticas, y muestra la posibilidad de que una política general supere a las políticas especializadas en una sola tarea.

\begin{figure}[H]
\centering
\includegraphics[width=0.96\textwidth]{ch24-05336.jpg}
\caption{RT-X / Open X-Embodiment: conjunto de datos de múltiples robots y el modelo RT-X. Fotograma de la pantalla proyectada, aprox. 01:28:56.}
\end{figure}

\begin{knowledgebox}{Lectura de la figura: el conjunto de datos es en sí mismo infraestructura}
La figura reúne una gran cantidad de robots y escenarios de tareas. Al leerla, conviene notar que el cuello de botella de VLA no es solo la estructura del modelo, sino también la estandarización de datos entre laboratorios, entre plataformas de hardware y entre tareas.
\end{knowledgebox}

\subsection{OpenVLA y HiRT: código abierto y procesamiento jerárquico de acciones}

OpenVLA se aproxima a una versión de código abierto de RT-2: conecta un VLM a la salida de acciones y publica como código abierto el modelo y el proceso de entrenamiento. HiRT, por su parte, señala que producir acciones directamente con un VLM presenta problemas de frecuencia de acción y de control fino, por lo que añade una policy/head dedicada a procesar la acción (Action): el VLM se encarga, a baja frecuencia, de la comprensión y la planificación, y la Action Policy ejecuta el control a alta frecuencia.

\begin{figure}[H]
\centering
\includegraphics[width=0.96\textwidth]{ch25-05531.jpg}
\caption{OpenVLA: modelo VLA de código abierto que conecta un VLM con la salida de acciones del robot. Fotograma de la pantalla proyectada, aprox. 01:32:11.}
\end{figure}

\begin{figure}[H]
\centering
\includegraphics[width=0.96\textwidth]{ch26-05765.jpg}
\caption{HiRT: mejora del control de robots mediante Robot Transformers jerárquicos. Fotograma de la pantalla proyectada, aprox. 01:36:05.}
\end{figure}

\begin{knowledgebox}{Lectura de la figura: el Action Head es un complemento importante para llevar VLA a la ingeniería}
En la figura de HiRT, la información del VLM de alto nivel se transmite a la Action Policy. Esto corresponde a un hecho de ingeniería: un VLM grande infiere con lentitud, el control de robots opera a alta frecuencia, y la salida de acciones necesita un módulo dedicado para manejar el control continuo y la retroalimentación visual local.
\end{knowledgebox}

\subsection{Figure Helix, pi0 y GR00T N1}

El siguiente tramo entra en la industria y en sistemas más cercanos a la conversión en producto. OpenVLA y HiRT ya dejaron al descubierto una tensión central: el VLM es bueno para la comprensión a baja frecuencia y la planificación semántica, pero un robot real necesita acciones (action) de alta frecuencia, continuas y estables. Lo que Figure Helix, pi0 y GR00T N1 tienen en común es que separan con mayor claridad el «módulo de comprensión» y el «módulo de acción», y luego los vuelven a conectar mediante el diseño del sistema.

Aunque Figure Helix no ha publicado un artículo, representa la arquitectura más reciente de la industria: el Sistema 2, en el nivel superior, funciona como un VLM preentrenado, y el Sistema 1, en el nivel inferior, funciona como control de acciones en tiempo real. pi0 introduce un flow model en VLA y usa un Action Expert para el control general de robots. NVIDIA GR00T N1, por su parte, pone el énfasis en un modelo fundacional abierto para robots humanoides y también adopta una estructura de VLM más procesamiento de acciones.

\begin{figure}[H]
\centering
\includegraphics[width=0.92\textwidth]{ch27-05946.jpg}
\caption{Figure AI Helix: esquema de la arquitectura de Sistema 2 y Sistema 1. Fotograma de la pantalla proyectada, aprox. 01:39:06.}
\end{figure}

\begin{figure}[H]
\centering
\includegraphics[width=0.92\textwidth]{ch28-06038.jpg}
\caption{pi0: Vision-Language-Action Flow Model for General Robot Control. Fotograma de la pantalla proyectada, aprox. 01:40:38.}
\end{figure}

\begin{figure}[H]
\centering
\includegraphics[width=0.94\textwidth]{ch29-06126.jpg}
\caption{NVIDIA GR00T N1: modelo fundacional abierto para robots humanoides de propósito general. Fotograma de la pantalla proyectada, aprox. 01:42:06.}
\end{figure}

\subsection{Resumen de la sección}

RT-2 transfiere el conocimiento de un VLM al control de robots; RT-X/Open X-Embodiment pone el énfasis en la escala de los datos y en el funcionamiento entre plataformas robóticas; OpenVLA ofrece una vía de código abierto, y HiRT/pi0/GR00T N1 comienzan a tratar con más seriedad el módulo de acción (Action). Esto indica que la dificultad central de VLA se está desplazando de «conectar el lenguaje y la visión» a «cómo producir acciones con alta frecuencia, de forma estable y con capacidad de generalización».
